\documentclass[11pt,a4paper]{report}

\usepackage[utf8]{inputenc}
\usepackage[T1]{fontenc}
\usepackage[french]{babel}
\usepackage{lmodern}
\usepackage{geometry}
\usepackage{xcolor}
\usepackage{graphicx}
\usepackage{booktabs}
\usepackage{multirow}
\usepackage{array}
\usepackage{longtable}
\usepackage{tabularx}
\usepackage{listings}
\usepackage{fancyhdr}
\usepackage{titlesec}
\usepackage{hyperref}
\usepackage{enumitem}
\usepackage{float}
\usepackage{caption}
\usepackage{microtype}

\geometry{top=2.2cm,bottom=2.2cm,left=2.4cm,right=2.4cm}
\definecolor{cmrpiblue}{RGB}{20,67,107}
\definecolor{cmrpilight}{RGB}{232,241,247}
\definecolor{codegreen}{RGB}{20,110,70}
\hypersetup{
  colorlinks=true,
  linkcolor=cmrpiblue,
  urlcolor=cmrpiblue,
  pdftitle={Rapport du Jalon 2 - Détecteur d'URL de phishing},
  pdfauthor={ALLADO Kossi Richard et ADJEVI Kodjo Marius-Ben}
}

\pagestyle{fancy}
\fancyhf{}
\fancyhead[L]{Jalon 2 -- Projet de stage PFA}
\fancyhead[R]{Détection d'URL de phishing}
\fancyfoot[L]{ALLADO Kossi Richard -- ADJEVI Kodjo Marius-Ben}
\fancyfoot[R]{Page \thepage}
\renewcommand{\headrulewidth}{0.4pt}
\renewcommand{\footrulewidth}{0.4pt}

\titleformat{\chapter}[hang]
  {\normalfont\huge\bfseries\color{cmrpiblue}}
  {\thechapter.}{0.7em}{}
\titlespacing*{\chapter}{0pt}{-20pt}{18pt}
\titleformat{\section}
  {\normalfont\Large\bfseries\color{cmrpiblue}}
  {\thesection}{0.7em}{}

\lstset{
  language=Python,
  basicstyle=\ttfamily\small,
  keywordstyle=\color{cmrpiblue}\bfseries,
  commentstyle=\color{codegreen},
  stringstyle=\color{purple},
  numbers=left,
  numberstyle=\tiny,
  stepnumber=1,
  numbersep=8pt,
  showstringspaces=false,
  breaklines=true,
  frame=single,
  backgroundcolor=\color{cmrpilight},
  captionpos=b
}

\newcommand{\metric}[2]{%
  \begin{minipage}[t]{0.22\textwidth}
    \centering
    \colorbox{cmrpilight}{\parbox[c][1.6cm][c]{0.92\linewidth}{%
      \centering\Large\bfseries\color{cmrpiblue}#1\\[-1mm]
      \normalsize\color{black}#2}}
  \end{minipage}
}

\begin{document}

\begin{titlepage}
  \centering
  {\large\bfseries Centre Marocain de Recherches Polytechnique et d'Innovation\par}
  \vspace{0.25cm}
  {\large Département / Service : Cybersécurité IA\par}
  \vspace{1.3cm}
  \rule{\textwidth}{1.2pt}
  \vspace{0.5cm}

  {\Large\bfseries PROJET DE STAGE PFA -- ÉTÉ 2026\par}
  \vspace{0.5cm}
  {\LARGE\bfseries\color{cmrpiblue}
  Système de détection et d'alerte précoce\\
  des cybermenaces ciblant les PME\par}

  \vspace{0.5cm}
  \rule{\textwidth}{1.2pt}
  \vspace{1.3cm}

  {\Large\bfseries LIVRABLE DU DEUXIÈME JALON\par}
  \vspace{0.7cm}
  {\LARGE\bfseries
  Construction d'un premier détecteur\\
  d'URL de phishing\par}
  \vspace{0.4cm}
  {\large Feature Engineering, automatisation des données\\
  et régression logistique\par}

  \vfill

  \begin{tabularx}{\textwidth}{X X}
    \textbf{Réalisé par :} & \textbf{Encadré par :}\\[0.25cm]
    ALLADO Kossi Richard & Pr. Youssef Bentaleb\\
    ADJEVI Kodjo Marius-Ben & Co-encadrante : Pr. Mehdia Ajana
  \end{tabularx}

  \vfill
  {\large\bfseries Période du deuxième jalon\par}
  {\large Du 1\textsuperscript{er} au 15 août 2026\par}
  \vspace{0.7cm}
  {\large Année universitaire 2025--2026\par}
\end{titlepage}

\pagenumbering{roman}
\tableofcontents
\clearpage
\pagenumbering{arabic}

\chapter{Introduction}

Le projet de stage PFA intitulé \og Système de détection et d'alerte précoce
des cybermenaces ciblant les PME \fg{} vise à construire progressivement un
prototype capable de collecter, détecter et signaler des menaces liées au
phishing. Le premier jalon a permis de comprendre le phishing, de présenter
les URL malveillantes et de développer un script de collecte depuis la
plateforme publique URLhaus.

Le deuxième jalon transforme ces données brutes en informations exploitables
par un algorithme d'apprentissage automatique. L'objectif principal consiste
à construire un premier détecteur simple, compréhensible et reproductible.
Pour cela, nous avons préparé un jeu de données étiqueté, extrait des
caractéristiques numériques à partir des URL, entraîné une régression
logistique et évalué ses performances sur un ensemble de test indépendant.

À la suite du retour de l'encadrant, une attention particulière a également
été portée à deux améliorations :

\begin{itemize}
  \item enrichir les URL au moyen d'une phase de \emph{Feature Engineering} ;
  \item automatiser la collecte, le nettoyage, la mise à jour et
        l'historisation des données.
\end{itemize}

Ce rapport décrit la démarche, les choix réalisés, les résultats obtenus et
les limites du premier modèle. Le code source associé est disponible dans le
dépôt GitHub du projet :
\url{https://github.com/ALLAKORI/STAGE_PFA_CMRPI}.

\chapter{Objectifs et démarche du Jalon 2}

\section{Objectifs}

Les objectifs opérationnels du deuxième jalon sont les suivants :

\begin{enumerate}
  \item récupérer un jeu de données d'URL déjà étiquetées ;
  \item transformer les URL en variables numériques pertinentes ;
  \item séparer les données d'entraînement et de test ;
  \item entraîner un modèle simple avec \texttt{scikit-learn} ;
  \item mesurer les bonnes détections et les erreurs ;
  \item sauvegarder le modèle afin de prédire de nouvelles URL ;
  \item rendre la collecte de données plus robuste et reproductible.
\end{enumerate}

\section{Chaîne de traitement}

La chaîne mise en place suit les étapes ci-dessous :

\begin{center}
\fbox{\parbox{0.9\textwidth}{
\centering
Collecte ou chargement des URL
$\longrightarrow$ nettoyage
$\longrightarrow$ extraction des caractéristiques
$\longrightarrow$ séparation entraînement/test
$\longrightarrow$ standardisation
$\longrightarrow$ entraînement
$\longrightarrow$ évaluation et sauvegarde
}}
\end{center}

Cette organisation sépare les responsabilités entre plusieurs scripts. Le
fichier \texttt{Download\_urls.py} gère la collecte, le fichier
\texttt{feature\_engineering.py} transforme une URL, le fichier
\texttt{train\_model.py} entraîne le détecteur et le fichier
\texttt{predict.py} applique le modèle sauvegardé à une nouvelle URL.

\chapter{Données utilisées}

\section{Jeu de données étiqueté}

L'entraînement s'appuie sur le jeu de données \emph{PhiUSIIL Phishing URL
Dataset}. Celui-ci contient des URL légitimes et des URL de phishing,
accompagnées d'une étiquette de classe. Après préparation et extraction des
variables utiles, le fichier \texttt{data/processed/dataset\_features.csv} contient
235\,795 observations et 22 colonnes : 21 caractéristiques et une étiquette.

\begin{table}[H]
\centering
\caption{Répartition des classes dans le jeu préparé}
\begin{tabular}{lrr}
\toprule
\textbf{Classe} & \textbf{Signification} & \textbf{Nombre d'URL}\\
\midrule
0 & Phishing & 100\,945\\
1 & Légitime & 134\,850\\
\midrule
\multicolumn{2}{r}{\textbf{Total}} & \textbf{235\,795}\\
\bottomrule
\end{tabular}
\end{table}

Les deux classes sont suffisamment représentées pour effectuer un premier
apprentissage. La légère différence entre leurs effectifs est prise en compte
lors de la séparation grâce à l'option \texttt{stratify=y}, qui conserve des
proportions comparables dans les ensembles d'entraînement et de test.

\section{Différence entre URLhaus et le jeu d'entraînement}

URLhaus fournit des indicateurs réels et récents associés principalement à la
distribution de logiciels malveillants. Cette source est utile pour alimenter
et actualiser la base locale, mais elle ne fournit pas à elle seule un
échantillon équilibré d'URL légitimes et de phishing adapté à l'apprentissage
supervisé.

Le jeu PhiUSIIL est donc utilisé pour entraîner et évaluer le premier modèle,
tandis qu'URLhaus reste la source de Threat Intelligence destinée à la
collecte continue. Cette distinction évite de considérer toutes les données
comme interchangeables et clarifie le rôle de chaque source.

\chapter{Extraction des caractéristiques}

\section{Principe du Feature Engineering}

Un modèle de régression logistique ne peut pas interpréter directement une
chaîne de caractères telle qu'une URL. Le \emph{Feature Engineering} consiste
à transformer cette chaîne en valeurs numériques décrivant sa structure.
L'objectif n'est pas d'ouvrir le site, mais d'analyser uniquement l'adresse.

Cette analyse statique réduit le risque lié à la manipulation d'URL
malveillantes et permet une prédiction rapide. Cependant, elle ne renseigne
pas sur le contenu réel de la page, son certificat, son ancienneté ou son
comportement après redirection.

\section{Caractéristiques retenues}

\begin{longtable}{p{0.34\textwidth}p{0.56\textwidth}}
\caption{Les 21 caractéristiques extraites d'une URL}\\
\toprule
\textbf{Caractéristique} & \textbf{Description}\\
\midrule
\endfirsthead
\toprule
\textbf{Caractéristique} & \textbf{Description}\\
\midrule
\endhead
\texttt{url\_length} & Longueur totale de l'URL.\\
\texttt{domain\_length} & Longueur du nom d'hôte.\\
\texttt{path\_length} & Longueur du chemin après le domaine.\\
\texttt{query\_length} & Longueur de la chaîne de paramètres.\\
\texttt{number\_of\_subdomains} & Nombre approximatif de sous-domaines.\\
\texttt{has\_ip\_address} & Indique si une adresse IP remplace le domaine.\\
\texttt{uses\_https} & Indique la présence du protocole HTTPS.\\
\texttt{number\_of\_letters} & Nombre de lettres.\\
\texttt{number\_of\_digits} & Nombre de chiffres.\\
\texttt{number\_of\_special\_characters} & Nombre de caractères non alphanumériques.\\
\texttt{letter\_ratio} & Proportion de lettres dans l'URL.\\
\texttt{digit\_ratio} & Proportion de chiffres dans l'URL.\\
\texttt{number\_of\_dots} & Nombre de points.\\
\texttt{number\_of\_hyphens} & Nombre de tirets.\\
\texttt{number\_of\_underscores} & Nombre de tirets bas.\\
\texttt{number\_of\_slashes} & Nombre de barres obliques.\\
\texttt{number\_of\_question\_marks} & Nombre de points d'interrogation.\\
\texttt{number\_of\_equal\_signs} & Nombre de signes égal.\\
\texttt{number\_of\_ampersands} & Nombre d'esperluettes.\\
\texttt{number\_of\_at\_symbols} & Nombre d'arobases.\\
\texttt{suspicious\_keyword\_count} & Nombre de mots suspects tels que
\texttt{login}, \texttt{verify}, \texttt{account} ou \texttt{password}.\\
\bottomrule
\end{longtable}

\section{Extrait de l'implémentation}

\begin{lstlisting}[caption={Extrait de la fonction d'extraction}]
def extract_features(url):
    original_url = str(url).strip()
    normalized_url = normalize_url(original_url)
    parsed_url = urlparse(normalized_url)

    hostname = parsed_url.hostname or ""
    path = parsed_url.path or ""
    query = parsed_url.query or ""

    return {
        "url_length": len(original_url),
        "domain_length": len(hostname),
        "path_length": len(path),
        "query_length": len(query),
        "number_of_subdomains": count_subdomains(hostname),
        "has_ip_address": contains_ip_address(hostname),
        "uses_https": int(
            original_url.lower().startswith("https://")
        ),
        # autres caracteristiques...
    }
\end{lstlisting}

La fonction normalise d'abord l'adresse en ajoutant un protocole lorsqu'il est
absent. La fonction \texttt{urlparse()} sépare ensuite le domaine, le chemin et
les paramètres. L'adresse IP est vérifiée avec le module standard
\texttt{ipaddress}, sans effectuer de connexion vers l'URL.

\chapter{Automatisation de la collecte}

\section{Améliorations apportées}

Le premier script téléchargeait les données URLhaus et remplaçait le fichier
local. Pour le deuxième jalon, le processus a été enrichi afin de :

\begin{itemize}
  \item vérifier les erreurs HTTP et limiter l'attente à 30 secondes ;
  \item nettoyer les commentaires présents dans le CSV d'URLhaus ;
  \item normaliser les URL et supprimer les doublons ;
  \item comparer les nouvelles données à la base locale ;
  \item ajouter uniquement les URL encore inconnues ;
  \item conserver un historique daté des nouvelles collectes ;
  \item calculer un hash SHA-256 du contenu téléchargé ;
  \item éviter un retraitement lorsque la source n'a pas changé ;
  \item mémoriser la date de la dernière mise à jour.
\end{itemize}

\section{Détection d'un changement}

\begin{lstlisting}[caption={Vérification du contenu téléchargé}]
contenu_brut = response.text

hash_actuel = hashlib.sha256(
    contenu_brut.encode("utf-8")
).hexdigest()

etat = charger_etat()

if etat.get("hash") == hash_actuel:
    print("Aucun changement detecte.")
    return
\end{lstlisting}

Le hash représente une empreinte du fichier. Si l'empreinte actuelle est
identique à celle de la collecte précédente, le programme s'arrête sans
réécrire inutilement les données.

\section{Mise à jour et historique}

Les URL déjà présentes sont réunies dans un ensemble, puis les nouvelles
lignes sont sélectionnées par différence. Après fusion, les doublons sont
supprimés. Un fichier daté est créé dans
\texttt{data/collected/urlhaus/history} afin de conserver uniquement les ajouts de
chaque exécution.

Cette organisation améliore la robustesse, la maintenabilité et la
reproductibilité demandées par l'encadrant. Le script peut ensuite être lancé
périodiquement à l'aide du Planificateur de tâches Windows ou d'une tâche
\texttt{cron} sous Linux. La programmation automatique par le système
d'exploitation sera intégrée et documentée lors de l'assemblage final.

\chapter{Construction du modèle}

\section{Choix de la régression logistique}

La régression logistique est un algorithme supervisé simple destiné à la
classification. Malgré son nom, elle permet ici de choisir entre deux classes :
phishing ou légitime. Elle calcule une probabilité à partir d'une combinaison
pondérée des caractéristiques.

Ce modèle a été retenu parce qu'il est rapide, accessible, adapté à une
première classification binaire et plus facile à expliquer qu'un modèle
complexe. Il fournit également une méthode \texttt{predict\_proba()} utile
pour afficher un niveau de confiance.

\section{Séparation entraînement/test}

Le jeu de données est divisé en deux parties :

\begin{itemize}
  \item 80\,\% des observations servent à entraîner le modèle ;
  \item 20\,\% servent uniquement à mesurer ses performances.
\end{itemize}

\begin{lstlisting}[caption={Séparation stratifiée des données}]
X_train, X_test, y_train, y_test = train_test_split(
    X,
    y,
    test_size=0.2,
    random_state=42,
    stratify=y,
)
\end{lstlisting}

Le paramètre \texttt{random\_state=42} rend la séparation reproductible.
L'ensemble de test contient 47\,159 URL : 20\,189 de phishing et 26\,970
légitimes. Ces URL ne sont pas utilisées pendant l'apprentissage.

\section{Standardisation}

Les variables n'ont pas toutes la même échelle. Par exemple, une longueur
d'URL peut atteindre plusieurs dizaines ou centaines, alors que
\texttt{uses\_https} vaut uniquement 0 ou 1. La standardisation centre et
réduit les variables pour faciliter l'optimisation du modèle.

\begin{lstlisting}[caption={Standardisation sans fuite de données}]
scaler = StandardScaler()

X_train_scaled = scaler.fit_transform(X_train)
X_test_scaled = scaler.transform(X_test)
\end{lstlisting}

Le scaler est ajusté uniquement sur les données d'entraînement. L'ensemble de
test est transformé avec les paramètres déjà appris, ce qui évite de transmettre
au modèle des informations provenant du test.

\section{Entraînement et sauvegarde}

\begin{lstlisting}[caption={Entraînement de la régression logistique}]
model = LogisticRegression(
    max_iter=1000,
    random_state=42,
)

model.fit(X_train_scaled, y_train)

joblib.dump(model, "models/logistic_regression_model.pkl")
joblib.dump(scaler, "models/scaler.pkl")
\end{lstlisting}

Le modèle et le scaler sont enregistrés avec \texttt{joblib}. Une nouvelle
prédiction peut ainsi être réalisée sans recommencer l'entraînement.

\chapter{Évaluation et résultats}

\section{Métriques utilisées}

L'\emph{accuracy} mesure la proportion totale de prédictions correctes. La
précision indique, parmi les URL annoncées comme phishing, la proportion
réellement associée au phishing. Le rappel indique la proportion d'URL de
phishing effectivement détectées. Le F1-score fournit un équilibre entre
précision et rappel.

Dans notre codage, la classe 0 correspond au phishing et la classe 1 à une URL
légitime. Le script évalue explicitement la classe 0 pour les métriques
principales, car la détection du phishing est l'objectif prioritaire. Le rapport
de classification détaille ensuite les résultats pour les deux classes.

\section{Résultats obtenus}

\begin{center}
\metric{99,32\,\%}{Accuracy}\hfill
\metric{99,91\,\%}{Précision phishing}\hfill
\metric{98,49\,\%}{Rappel phishing}\hfill
\metric{99,20\,\%}{F1-score phishing}
\end{center}

\vspace{0.6cm}

\begin{table}[H]
\centering
\caption{Rapport de classification sur l'ensemble de test}
\begin{tabular}{lrrrr}
\toprule
\textbf{Classe} & \textbf{Précision} & \textbf{Rappel} &
\textbf{F1-score} & \textbf{Effectif}\\
\midrule
Phishing (0) & 1,00 & 0,98 & 0,99 & 20\,189\\
Légitime (1) & 0,99 & 1,00 & 0,99 & 26\,970\\
\midrule
Accuracy & \multicolumn{3}{c}{0,99} & 47\,159\\
Moyenne macro & 0,99 & 0,99 & 0,99 & 47\,159\\
Moyenne pondérée & 0,99 & 0,99 & 0,99 & 47\,159\\
\bottomrule
\end{tabular}
\end{table}

\section{Matrice de confusion}

\begin{table}[H]
\centering
\caption{Matrice de confusion}
\begin{tabular}{llrr}
\toprule
& & \multicolumn{2}{c}{\textbf{Classe prédite}}\\
& & Phishing (0) & Légitime (1)\\
\midrule
\multirow{2}{*}{\textbf{Classe réelle}} & Phishing (0) & 19\,885 & 304\\
& Légitime (1) & 18 & 26\,952\\
\bottomrule
\end{tabular}
\end{table}

\noindent Les lignes représentent les classes réelles et les colonnes les
classes prédites.

Sur 47\,159 URL de test :

\begin{itemize}
  \item 19\,885 URL de phishing sont correctement détectées ;
  \item 26\,952 URL légitimes sont correctement reconnues ;
  \item 304 URL de phishing sont classées à tort comme légitimes ;
  \item 18 URL légitimes sont classées à tort comme phishing.
\end{itemize}

L'erreur la plus importante pour la sécurité correspond aux 304 URL de
phishing classées comme légitimes, car elles pourraient ne pas déclencher
d'alerte. Le taux global est élevé, mais le prototype ne doit donc pas être
considéré comme une garantie absolue.

\section{Interprétation}

Les résultats montrent que les caractéristiques lexicales et structurelles
permettent de distinguer efficacement les classes dans ce jeu de données.
Cependant, une performance élevée sur un ensemble provenant de la même source
ne prouve pas que le modèle conservera exactement la même efficacité sur des
URL réelles, nouvelles ou issues d'une autre distribution.

Il sera donc nécessaire de tester le détecteur sur des données plus récentes,
de vérifier l'absence de doublons entre les ensembles et de surveiller
l'évolution des performances dans le temps.

\chapter{Prédiction d'une nouvelle URL}

Le script \texttt{predict.py} charge le modèle et le scaler enregistrés. Il
demande une URL à l'utilisateur, extrait les 21 caractéristiques dans le même
ordre que pendant l'entraînement, applique la standardisation, puis affiche la
classe et la probabilité associée.

\begin{lstlisting}[caption={Principe de la prédiction interactive}]
url = input("Entrez une URL a analyser : ").strip()

features = extract_features(url)
features_df = pd.DataFrame([features])
features_scaled = scaler.transform(features_df)

prediction = model.predict(features_scaled)[0]
probabilities = model.predict_proba(features_scaled)[0]
\end{lstlisting}

Si la prédiction vaut 1, le résultat affiché est \og URL légitime \fg. Si elle
vaut 0, le résultat est \og URL phishing \fg. La confiance affichée correspond
à la probabilité calculée pour la classe prédite.

Cette utilisation reste une aide à la décision. L'utilisateur ne doit pas
ouvrir une URL suspecte pour confirmer le résultat. Une adresse classée
légitime peut être malveillante, notamment si elle utilise une technique qui
n'est pas représentée dans le jeu d'entraînement.

\chapter{Organisation et reproductibilité}

\subsection{Organisation du dépôt}

Le dépôt du projet est organisé de manière à séparer les différentes étapes de la chaîne de traitement, depuis la collecte des données jusqu'à l'entraînement du modèle et son utilisation. Cette organisation facilite la maintenance du projet ainsi que la reproductibilité des expériences.

\begin{lstlisting}[language={},basicstyle=\ttfamily\small]
project/
|--- data/
|   |--- raw/
|   |   `--- PhiUSIIL_Phishing_URL_Dataset.csv
|   |--- processed/
|   |   `--- dataset_features.csv
|   `--- collected/
|       `--- urlhaus/
|           |--- current_urls.csv
|           `--- history/
|
|--- docs/
|   |--- evaluation.md
|   `--- report/
|
|--- models/
|   |--- logistic_regression_model.pkl
|   `--- scaler.pkl
|
|--- Download_urls.py
|--- feature_engineering.py
|--- train_model.py
|--- predict.py
|--- prepare_phishing_url_features.ipynb
|--- requirements.txt
`--- README.md
\end{lstlisting}

Les principaux éléments du dépôt sont décrits ci-dessous :

\begin{itemize}
    \item \textbf{data/raw/} : contient le jeu de données d'origine utilisé pour l'entraînement des modèles.
    \item \textbf{data/processed/} : regroupe les jeux de données après extraction des caractéristiques.
    \item \textbf{data/collected/urlhaus/} : stocke les URL téléchargées depuis URLhaus ainsi que l'historique des différentes collectes.
    \item \textbf{docs/} : rassemble la documentation technique, les résultats d'évaluation et les sources du rapport.
    \item \textbf{models/} : contient les modèles entraînés ainsi que les objets nécessaires à la phase de prédiction (standardisation, sérialisation).
    \item \textbf{Download\_urls.py} : automatise le téléchargement et la mise à jour des URL provenant d'URLhaus.
    \item \textbf{feature\_engineering.py} : extrait les 21 caractéristiques utilisées par les modèles de classification.
    \item \textbf{train\_model.py} : réalise l'entraînement, l'évaluation et la sauvegarde des modèles.
    \item \textbf{predict.py} : permet d'effectuer la prédiction du caractère malveillant d'une URL à partir d'un modèle entraîné.
    \item \textbf{prepare\_phishing\_url\_features.ipynb} : notebook utilisé pour l'exploration des données et l'extraction des caractéristiques.
    \item \textbf{requirements.txt} : liste des dépendances Python nécessaires à l'exécution du projet.
    \item \textbf{README.md} : présente les instructions d'installation, d'utilisation et l'organisation générale du dépôt.
\end{itemize}

\section{Reproduction de l'expérience}

\begin{lstlisting}[language={},caption={Installation et exécution sous Windows}]
python -m venv venv
.\venv\Scripts\Activate.ps1
pip install -r requirements.txt
python train_model.py
python predict.py
\end{lstlisting}

Le dépôt Git permet de conserver l'historique des modifications et de
partager une version commune entre les deux stagiaires. L'environnement
virtuel et les caches ne sont pas publiés. Les dépendances nécessaires sont
déclarées dans \texttt{requirements.txt}.

\chapter{Limites et perspectives}

\section{Limites actuelles}

\begin{itemize}
  \item Le modèle analyse la forme de l'URL, mais pas le contenu de la page.
  \item HTTPS ne garantit pas qu'un site est légitime.
  \item Les mots-clés suspects peuvent aussi apparaître dans une URL légitime.
  \item Les performances peuvent diminuer sur des données plus récentes.
  \item La confiance du modèle n'est pas une preuve de sécurité.
  \item URLhaus concerne plus largement les URL diffusant des logiciels
        malveillants et ne représente pas uniquement le phishing.
  \item Le prototype ne doit pas ouvrir automatiquement les URL collectées.
\end{itemize}

\section{Préparation du Jalon 3}

Le troisième jalon assemblera les éléments dans un prototype complet. Les
travaux prévus sont :

\begin{enumerate}
  \item créer un tableau de bord Streamlit ;
  \item afficher les informations et le résultat de détection ;
  \item envoyer une alerte par e-mail lorsqu'une menace est détectée ;
  \item rédiger deux à trois fiches réflexes destinées aux PME ;
  \item tester la chaîne de bout en bout ;
  \item préparer la documentation et la présentation finale.
\end{enumerate}

Une attention particulière sera portée à la protection des identifiants de
messagerie. Ils ne devront pas être enregistrés dans le code ni publiés sur
GitHub. Ils seront chargés à partir de variables d'environnement ou d'un
fichier local exclu par \texttt{.gitignore}.

\chapter{Conclusion}

Ce deuxième jalon a permis de construire un premier détecteur d'URL de
phishing fondé sur une régression logistique. Les URL ont été transformées en
21 caractéristiques numériques décrivant leur longueur, leur domaine, leurs
caractères et certains motifs suspects.

Le jeu de 235\,795 observations a été séparé en ensembles d'entraînement et
de test de manière stratifiée. Sur les 47\,159 observations réservées au test,
le modèle obtient une accuracy de 99,32\,\% et un F1-score de 99,41\,\%.
L'analyse de la matrice de confusion montre 322 erreurs, dont 304 URL de
phishing classées comme légitimes.

Le processus de collecte URLhaus a également été amélioré avec la suppression
des doublons, la détection des changements, la mise à jour incrémentale et la
conservation d'un historique. Ces améliorations répondent aux recommandations
formulées après le premier jalon.

Le modèle constitue une base fonctionnelle pour le prototype, mais il conserve
des limites et ne doit pas être utilisé comme unique preuve de sécurité. Le
prochain jalon intégrera ce détecteur dans une interface Streamlit avec un
mécanisme d'alerte et des fiches réflexes destinées aux PME.

\chapter*{Références}
\addcontentsline{toc}{chapter}{Références}

\begin{itemize}
  \item URLhaus, abuse.ch : \url{https://urlhaus.abuse.ch}
  \item Documentation scikit-learn :
        \url{https://scikit-learn.org/stable/}
  \item Documentation pandas : \url{https://pandas.pydata.org/docs/}
  \item Documentation Python : \url{https://docs.python.org/3/}
  \item PhiUSIIL Phishing URL Dataset, UCI Machine Learning Repository :
        \url{https://archive.ics.uci.edu/dataset/967/phiusiil+phishing+url+dataset}
  \item Dépôt GitHub du projet :
        \url{https://github.com/ALLAKORI/STAGE_PFA_CMRPI}
\end{itemize}

\appendix
\chapter{Valeurs détaillées de l'exécution}

\begin{verbatim}
Accuracy : 0.993172035030429
Precision : 0.9888464925154095
Recall : 0.9993325917686318
F1-score : 0.9940618891306753

Matrice de confusion :
[[19885   304]
 [   18 26952]]
\end{verbatim}

\end{document}
